\chapter{Variabili aleatorie continue multiple}

Nel Capitolo~8 abbiamo descritto coppie di variabili discrete con una tabella, la pmf congiunta. Qui facciamo lo
stesso con le variabili continue: la tabella diventa una \textbf{superficie}, la densità congiunta, e le probabilità
diventano \textbf{volumi} sotto la superficie. Per calcolarli ci serve l'integrale doppio, che spieghiamo da zero.
Tutto il resto (marginali, indipendenza, densità condizionate, media condizionata, covarianza) si traduce dal discreto
con le regole del dizionario del capitolo precedente.

% ============================================================
\section{Coppie e \texorpdfstring{$m$}{m}-ple di variabili continue}
% ============================================================

In perfetta analogia con quanto fatto per le variabili discrete, una coppia di variabili continue è definita nella
forma
\[ (X, Y) : \omega \in \Omega \;\longmapsto\; \left( X(\omega), Y(\omega) \right) \in \mathcal{X} \times \mathcal{Y}
   \subseteq \R^2 , \]
dove $\mathcal{X}$ e $\mathcal{Y}$ sono gli alfabeti di $X$ e di $Y$. Analogamente, date tre variabili aleatorie (a
questo punto non importa più se tutte continue o meno) $X(\omega) \in \mathcal{X}$, $Y(\omega) \in \mathcal{Y}$,
$Z(\omega) \in \mathcal{Z}$, si ha la terna
\[ (X, Y, Z) : \omega \in \Omega \;\longmapsto\; \left( X(\omega), Y(\omega), Z(\omega) \right) \in \mathcal{X} \times
   \mathcal{Y} \times \mathcal{Z} \subseteq \R^3 , \]
e, date $m$ variabili aleatorie $X_i \in \mathcal{X}_i \subseteq \R$, la $m$-pla aleatoria
\[ (X_1, X_2, \ldots, X_m) : \omega \in \Omega \;\longmapsto\; \left( X_1(\omega), \ldots, X_m(\omega) \right) \in
   \mathcal{X}_1 \times \mathcal{X}_2 \times \cdots \times \mathcal{X}_m \subseteq \R^m . \]

Perché servono? Si pensi ad altezza e peso di una persona scelta a caso in una popolazione. Le funzioni $F_X(x)$ e
$F_Y(y)$ caratterizzano il comportamento statistico delle singole variabili, ma non quello \textbf{congiunto}: ci
aspettiamo che una persona molto alta sia anche abbastanza pesante, e questa influenza reciproca non può essere
descritta da due funzioni che riguardano ciascuna una sola delle due grandezze. Serve una caratterizzazione congiunta.

% ============================================================
\section{La densità congiunta}
% ============================================================

\dfn{Densità congiunta}{
  La \term{densità di probabilità congiunta} della coppia di variabili continue $(X, Y)$ è
  \[ f_{X,Y}(x, y) = \lim_{\Delta x \to 0}\, \lim_{\Delta y \to 0} \frac{P\left( \left\{ x - \frac{\Delta x}{2} \leq X
     \leq x + \frac{\Delta x}{2} \right\} \cap \left\{ y - \frac{\Delta y}{2} \leq Y \leq y + \frac{\Delta y}{2} \right\}
     \right)}{\Delta x\, \Delta y} = \lim_{\Delta x \to 0}\, \lim_{\Delta y \to 0} \frac{P_{X,Y}(x, y; \Delta x,
     \Delta y)}{\Delta x\, \Delta y} , \]
  per $(x, y) \in \mathcal{X} \times \mathcal{Y}$.
}

È la probabilità che il punto $(X, Y)$ cada in un rettangolino attorno a $(x, y)$, divisa per l'\textbf{area} del
rettangolino: una probabilità per unità di superficie. Se nel caso di una variabile la probabilità era uno strato di
sabbia lungo una retta, ora è uno strato di sabbia steso su un tavolo, più alto in certi punti e più basso in altri.
La densità congiunta è l'altezza dello strato, e la probabilità che il punto cada in una regione $C$ del tavolo è il
\textbf{volume} di sabbia sopra quella regione:
\[ P\left( (X, Y) \in C \right) = \iint_C f_{X,Y}(x, y) \dd x \dd y . \]

\begin{richiamo}{l'integrale doppio}
Sia $f(x, y) \geq 0$ una funzione di due variabili: il suo grafico è una superficie sopra il piano $(x, y)$. Il
\term{volume} compreso fra la superficie e una regione $C$ del piano si calcola come l'area nel caso di una variabile:
si divide il piano in tanti quadratini di lati $\Delta x$ e $\Delta y$, sopra ogni quadratino si costruisce una colonna
alta quanto la funzione, $f(x_i, y_j)$, e si sommano i volumi delle colonne. Quando i quadratini diventano
infinitamente piccoli, la somma diventa l'\term{integrale doppio}:
\[ \iint_C f(x, y) \dd x \dd y = \lim_{\Delta x, \Delta y \to 0} \sum_{(x_i, y_j) \in C} f(x_i, y_j)\, \Delta x\, \Delta y . \]
\textbf{Come si calcola: a fette.} Sopra un rettangolo $[a, b] \times [c, d]$ si può prima tagliare il volume a fette
sottili perpendicolari all'asse $x$: per $x$ fissato, la fetta ha area $\int_c^d f(x, y) \dd y$ (un integrale semplice,
in cui $x$ è una costante). Poi si sommano le fette, integrando lungo $x$:
\[ \iint_{[a,b] \times [c,d]} f(x, y) \dd x \dd y = \int_a^b \left( \int_c^d f(x, y) \dd y \right) \dd x = \int_c^d
   \left( \int_a^b f(x, y) \dd x \right) \dd y . \]
L'ordine non conta (è il \term{teorema di Fubini}): si possono tagliare le fette in un verso o nell'altro, come una somma
doppia si può fare per righe o per colonne (Capitolo~6).
\begin{center}
\begin{tikzpicture}
\begin{axis}[width=7.6cm, height=6.2cm, view={-50}{22}, xlabel=$x$, ylabel=$y$, zlabel={}, xmin=0, xmax=1, ymin=0,
  ymax=1, zmin=0, zmax=2.1, xtick={0,1}, ytick={0,1}, ztick={1,2}, tick label style={font=\scriptsize},
  label style={font=\small}]
  \addplot3[surf, shader=flat, draw=mspiric!60, fill=mspiric!12, domain=0:1, y domain=0:1, samples=9] {x+y};
  \addplot3[fill=evid, draw=mspiB!85!black, line width=0.8pt] coordinates {(0.5,0,0) (0.5,1,0) (0.5,1,1.5) (0.5,0,0.5) (0.5,0,0)};
\end{axis}
\end{tikzpicture}
\hspace{0.3cm}
\begin{minipage}[b]{6.4cm}\small
La superficie $z = x + y$ sopra il quadrato $[0, 1]^2$. La fetta gialla, a $x = \frac12$ fissato, ha area
$\int_0^1 \left( \frac12 + y \right) \dd y = 1$. In generale la fetta in $x$ ha area $\int_0^1 (x + y) \dd y = x +
\frac12$, e sommando le fette
\[ \int_0^1 \left( x + \tfrac12 \right) \dd x = \tfrac12 + \tfrac12 = 1 . \]
\end{minipage}
\end{center}
\textbf{Due casi utili.} Se la funzione è un prodotto, $f(x, y) = g(x)\, h(y)$, sopra un rettangolo l'integrale doppio è
il prodotto di due integrali semplici: $\iint g(x) h(y) \dd x \dd y = \int g(x) \dd x \cdot \int h(y) \dd y$ (nella fetta,
$g(x)$ è una costante che esce dall'integrale in $y$). Se invece la regione non è un rettangolo, gli estremi
dell'integrale interno dipendono dalla variabile esterna: per il triangolo $x \geq 0$, $y \geq 0$, $x + y \leq 1$, la
fetta in $x$ va da $y = 0$ a $y = 1 - x$, e
$\iint_{\text{triangolo}} f \dd x \dd y = \int_0^1 \left( \int_0^{1 - x} f(x, y) \dd y \right) \dd x$.
\end{richiamo}

\teorema[proprietacongiunta]{Proprietà della densità congiunta}{
  Le densità congiunte devono soddisfare vincoli simili a quelli delle densità di una sola variabile:
  \begin{enumerate}[label=\alph*)]
    \item $f_{X,Y}(x, y) \geq 0$ per ogni $(x, y)$;
    \item $f_{X,Y}$ è sommabile su $\R^2$ con integrale unitario:
          $\displaystyle \int_{-\infty}^{\infty} \int_{-\infty}^{\infty} f_{X,Y}(x, y) \dd x \dd y = P\left( (X, Y) \in
          \R^2 \right) = 1$.
  \end{enumerate}
}

\begin{dimostrazione}
\passo{} La densità è il limite di rapporti fra una probabilità e un'area, quindi non è negativa; il volume totale è la
probabilità dell'evento certo «il punto cade da qualche parte nel piano».
\end{dimostrazione}

\begin{esempio}{Due tempi di attesa indipendenti}
Due componenti si guastano dopo tempi $X$ e $Y$ indipendenti, entrambi esponenziali di parametro 1. Come vedremo fra poco,
per variabili indipendenti la densità congiunta è il prodotto delle densità:
\[ f_{X,Y}(x, y) = e^{-x}\, e^{-y} = e^{-(x + y)} \qquad x, y \geq 0 . \]
Qual è la probabilità che il tempo totale $X + Y$ non superi 1? È il volume sopra il triangolo $x, y \geq 0$,
$x + y \leq 1$:
\[
  P(X + Y \leq 1) = \int_0^1 \left( \int_0^{1 - x} e^{-x} e^{-y} \dd y \right) \dd x = \int_0^1 e^{-x} \left( 1 -
  e^{-(1 - x)} \right) \dd x = \int_0^1 \left( e^{-x} - e^{-1} \right) \dd x = \left( 1 - e^{-1} \right) - e^{-1} = 1 -
  \frac{2}{e} \approx 0{,}264 .
\]
Nella fetta in $x$ il fattore $e^{-x}$ è costante ed esce dall'integrale; poi $e^{-x} e^{-(1-x)} = e^{-1}$.
\begin{center}
\begin{tikzpicture}
\begin{axis}[width=8.4cm, height=6.6cm, view={125}{30}, xlabel=$x$, ylabel=$y$, zlabel={$f_{X,Y}$},
  xmin=0, xmax=2.5, ymin=0, ymax=2.5, zmin=0, zmax=1, xtick={0,1,2}, ytick={0,1,2}, ztick={0.5,1},
  zticklabels={$0{,}5$,$1$}, tick label style={font=\scriptsize}, label style={font=\small}, colormap/viridis]
  \addplot3[surf, shader=flat, opacity=0.55, draw=black!25, domain=0:2.5, y domain=0:2.5, samples=21] {exp(-x-y)};
  \addplot3[fill=evid, draw=mspiB!85!black, line width=0.8pt] coordinates {(0,0,0) (1,0,0) (0,1,0) (0,0,0)};
  \addplot3[mspiB!85!black, line width=0.8pt, domain=0:1, samples=20, samples y=0] ({x},{1-x},{exp(-1)});
  \addplot3[mspiB!85!black, line width=0.8pt, domain=0:1, samples=20, samples y=0] ({x},{0},{exp(-x)});
  \addplot3[mspiB!85!black, line width=0.8pt, domain=0:1, samples=20, samples y=0] ({0},{x},{exp(-x)});
  \addplot3[mspiB!85!black, dashed] coordinates {(1,0,0) (1,0,0.3679)};
  \addplot3[mspiB!85!black, dashed] coordinates {(0,1,0) (0,1,0.3679)};
  \addplot3[mspiB!85!black, dashed] coordinates {(0,0,0) (0,0,1)};
\end{axis}
\end{tikzpicture}
\end{center}
La probabilità è il volume del «solido» che sta sopra il triangolo giallo e sotto la superficie $e^{-(x+y)}$.
\end{esempio}

% ============================================================
\section{Marginali, indipendenza e densità condizionate}
% ============================================================

La densità congiunta $f_{X,Y}$ condivide con la pmf congiunta $p_{X,Y}$ le stesse proprietà fondamentali.

\teorema[marginalic]{Marginalizzazione e indipendenza}{
  \begin{enumerate}[label=\alph*)]
    \item \textbf{Marginalizzazione}:
          \[ f_X(x) = \int_{\R} f_{X,Y}(x, y) \dd y , \qquad\qquad f_Y(y) = \int_{\R} f_{X,Y}(x, y) \dd x , \]
          per cui caratterizzare congiuntamente $(X, Y)$ significa anche caratterizzarle marginalmente, mentre il
          viceversa non è necessariamente vero.
    \item \textbf{Indipendenza statistica}: due variabili aleatorie sono indipendenti se e solo se
          \[ f_{X,Y}(x, y) = f_X(x)\, f_Y(y) \quad\iff\quad F_{X,Y}(x, y) = F_X(x)\, F_Y(y) . \]
    \item Più in generale, $X_1, \ldots, X_m$ con densità $f_{X_i}$ sono indipendenti se e solo se
          \[ f_{X_1, \ldots, X_m}(x_1, \ldots, x_m) = \prod_{i=1}^{m} f_{X_i}(x_i) \qquad \forall (x_1, \ldots, x_m) \in
             \R^m . \]
  \end{enumerate}
}

\begin{dimostrazione}[della marginalizzazione]
\passo{Una striscia} L'evento $\left\{ x - \frac{\Delta x}{2} \leq X \leq x + \frac{\Delta x}{2} \right\}$ non pone
condizioni su $Y$: nel piano è una \textbf{striscia verticale} larga $\Delta x$, illimitata in alto e in basso.
\passo{Il volume della striscia} La sua probabilità è il volume sopra la striscia. Per $\Delta x$ piccolo, ogni fetta della
striscia ha circa la stessa sezione, quella in $x$, cioè $\int_{\R} f_{X,Y}(x, y) \dd y$; quindi il volume è circa
\[ P\left( x - \tfrac{\Delta x}{2} \leq X \leq x + \tfrac{\Delta x}{2} \right) \simeq \Delta x \int_{\R} f_{X,Y}(x, y)
   \dd y . \]
\passo{La densità} Dividendo per $\Delta x$ e mandandolo a zero si ottiene $f_X(x)$. Per $f_Y$ si usa una striscia
orizzontale.
\end{dimostrazione}

\begin{figure}[H]
\centering
\begin{tikzpicture}[x=1.6cm, y=1.6cm]
  \foreach \i in {0,...,3} \foreach \j in {0,...,3} {
    \pgfmathsetmacro{\c}{12+18*(\i+\j)/6}
    \fill[mspiA!\c] (\i*0.25,\j*0.25) rectangle ++(0.25,0.25);
  }
  \fill[evid, opacity=0.85] (0.55,0) rectangle (0.65,1);
  \draw[mspigray] (0,0) rectangle (1,1);
  \draw[-{Stealth}] (-0.1,0) -- (1.3,0) node[below, font=\small] {$x$};
  \draw[-{Stealth}] (0,-0.1) -- (0,1.3) node[left, font=\small] {$y$};
  \node[font=\scriptsize, below] at (0.6,0) {$x$};
  \draw[decorate, decoration={brace, amplitude=3pt}] (0.55,1.03) -- (0.65,1.03) node[midway, above=3pt,
    font=\scriptsize] {$\Delta x$};
  \node[font=\scriptsize, below] at (1,0) {1}; \node[font=\scriptsize, left] at (0,1) {1};
  \node[font=\small, anchor=west, text width=7.6cm] at (1.5,0.5)
    {La densità $x + y$ sul quadrato (più scura dove è più alta). La striscia gialla è l'evento «$X$ cade vicino a
     $x$», qualunque sia $Y$: il suo volume diviso $\Delta x$ è la densità marginale $f_X(x) = \int_0^1 (x + y) \dd y =
     x + \frac12$.};
\end{tikzpicture}
\caption{Marginalizzare vuol dire «schiacciare» la densità congiunta su un asse, sommando lungo l'altro.}
\end{figure}

\subsection{Le densità condizionate}

Vorremmo la densità di $X$ sapendo che $Y = y$. Ma l'evento $\{Y = y\}$ ha probabilità zero, e non ci si può dividere
per zero. Si condiziona allora all'evento $\left\{ y - \frac{\Delta y}{2} \leq Y \leq y + \frac{\Delta y}{2}
\right\}$, che ha probabilità positiva, e solo alla fine si stringe l'intervallo:
\[
  P\left( x - \tfrac{\Delta x}{2} \leq X \leq x + \tfrac{\Delta x}{2} \;\Big|\; y - \tfrac{\Delta y}{2} \leq Y \leq y +
  \tfrac{\Delta y}{2} \right) = \frac{P_{X,Y}(x, y; \Delta x, \Delta y)}{P_Y(y; \Delta y)} .
\]
Dividendo per $\Delta x$ e mandando $\Delta x \to 0$ si ottiene la densità di $X$ condizionata all'evento «$Y$ vicino a
$y$». Facendo poi tendere anche $\Delta y$ a zero, e moltiplicando e dividendo per $\Delta y$:
\[
  f_{X \mid Y}(x \mid y) = \lim_{\Delta y \to 0}\, \lim_{\Delta x \to 0} \frac{P_{X,Y}(x, y; \Delta x, \Delta y)}{\Delta x\,
  P_Y(y; \Delta y)} = \frac{\displaystyle \lim_{\Delta y \to 0}\, \lim_{\Delta x \to 0} \frac{P_{X,Y}(x, y; \Delta x,
  \Delta y)}{\Delta x\, \Delta y}}{\displaystyle \lim_{\Delta y \to 0} \frac{P_Y(y; \Delta y)}{\Delta y}} =
  \frac{f_{X,Y}(x, y)}{f_Y(y)} .
\]

\dfn{Densità condizionata}{
  La \term{densità di $X$ condizionata a $Y = y$} (per $f_Y(y) > 0$) è
  \[ f_{X \mid Y}(x \mid y) = \frac{f_{X,Y}(x, y)}{f_Y(y)} . \]
}

Come c'era da attendersi, l'espressione riproduce l'analoga definizione per la pmf condizionata, che d'altronde è essa
stessa una densità (discreta). Di conseguenza tutte le proprietà delle pmf condizionate si estendono alle densità
condizionate.

\teorema[proprietacondc]{Proprietà delle densità condizionate}{
  \begin{enumerate}[label=\alph*)]
    \item Per $y$ fissato, al variare di $x$, $f_{X \mid Y}(x \mid y)$ è una densità:
          $f_{X \mid Y}(x \mid y) \geq 0$ e $\int_{\R} f_{X \mid Y}(x \mid y) \dd x = 1$.
    \item \textbf{Probabilità totale per le densità}:
          \[ f_X(x) = \int_{\R} f_{X,Y}(x, y) \dd y = \int_{\R} f_{X \mid Y}(x \mid y)\, f_Y(y) \dd y , \qquad
             f_Y(y) = \int_{\R} f_{Y \mid X}(y \mid x)\, f_X(x) \dd x . \]
    \item \textbf{Probabilità composta e Bayes}:
          \[ f_{X,Y}(x, y) = f_Y(y)\, f_{X \mid Y}(x \mid y) = f_X(x)\, f_{Y \mid X}(y \mid x) , \qquad
             f_{X \mid Y}(x \mid y) = \frac{f_X(x)\, f_{Y \mid X}(y \mid x)}{f_Y(y)} . \]
  \end{enumerate}
}

\begin{dimostrazione}
\passo{a)} È un rapporto di quantità non negative; integrando in $x$, il numeratore dà la marginale $f_Y(y)$, che si
semplifica con il denominatore.
\passo{b) e c)} Dalla definizione, $f_{X,Y} = f_{X \mid Y}\, f_Y$ (e, scambiando i ruoli, $= f_{Y \mid X}\, f_X$); sostituendo
nella marginalizzazione si ottiene b), uguagliando le due scritture e dividendo per $f_Y$ si ottiene Bayes.
\end{dimostrazione}

\begin{esempio}{Una densità congiunta calcolata fino in fondo}
Sia $f_{X,Y}(x, y) = x + y$ sul quadrato $[0, 1]^2$ (e 0 fuori). È una densità: è positiva e il suo integrale vale 1
(lo abbiamo calcolato nel richiamo).
\begin{itemize}
  \item \textbf{Marginali}: $f_X(x) = \int_0^1 (x + y) \dd y = x + \frac12$ e, per simmetria, $f_Y(y) = y + \frac12$,
        per $x, y \in [0, 1]$.
  \item \textbf{Indipendenza}: $f_X(x)\, f_Y(y) = \left(x + \frac12\right)\left(y + \frac12\right) \neq x + y$ (in
        $(0, 0)$ vale $\frac14$ contro 0): le variabili non sono indipendenti.
  \item \textbf{Condizionata}: $f_{X \mid Y}(x \mid y) = \dfrac{x + y}{y + \frac12}$, per $x \in [0, 1]$. Per esempio
        con $y = 0$ è $2x$ (sapendo che $Y$ è piccola, $X$ tende a essere grande), con $y = 1$ è $\frac{x + 1}{1{,}5}$,
        quasi piatta.
\end{itemize}
\begin{center}
\begin{tikzpicture}
\begin{axis}[quaderno, width=7.6cm, height=4.8cm, xmin=0, xmax=1.25, ymin=0, ymax=2.3, xtick={0.5,1}, ytick={1,2},
  xlabel=$x$, legend style={font=\scriptsize, at={(0.02,0.98)}, anchor=north west}]
  \addplot[mspiA, domain=0:1] {2*x};
  \addplot[mspiC, domain=0:1] {x+0.5};
  \addplot[mspiB, domain=0:1] {(x+1)/1.5};
  \legend{$f_{X \mid Y}(x \mid 0)$, $f_{X \mid Y}(x \mid \frac12)$, $f_{X \mid Y}(x \mid 1)$}
\end{axis}
\end{tikzpicture}
\end{center}
Ogni curva è una «fetta» della densità congiunta a $y$ fissato, divisa per la sua area $f_Y(y)$ in modo da avere area 1.
\end{esempio}

% ============================================================
\section{Medie, media condizionata e covarianza}
% ============================================================

Come nel caso discreto valgono le seguenti estensioni.
\begin{itemize}
  \item Se $Z = g(X, Y)$, allora
        \[ E[Z] = \iint_{\R^2} g(x, y)\, f_{X,Y}(x, y) \dd x \dd y . \]
  \item \textbf{Linearità della media}: $E\left[ \sum_{i=1}^{m} a_i X_i \right] = \sum_{i=1}^{m} a_i\, E[X_i]$.
\end{itemize}

\teorema[mediacondc]{Teorema della media condizionata}{
  \[ E[g(X, Y)] = E_Y\Big[ E_{X \mid Y}\big[ g(X, Y) \mid Y \big] \Big] , \qquad\text{dove}\qquad E[g(X, Y) \mid Y = y] =
     \int_{\R} g(x, y)\, f_{X \mid Y}(x \mid y) \dd x . \]
}

\begin{dimostrazione}
\passo{Una funzione di $y$} $E[g(X, Y) \mid Y = y]$ è una funzione di $y$, chiamiamola $h(y)$. Se consideriamo la variabile
aleatoria $h(Y)$ e ne calcoliamo la media statistica:
\[ E[h(Y)] = \int_{\R} h(y)\, f_Y(y) \dd y = \int_{\R} \left( \int_{\R} g(x, y)\, f_{X \mid Y}(x \mid y) \dd x \right)
   f_Y(y) \dd y . \]
\passo{Si riconosce la congiunta} $f_Y(y)$ non dipende da $x$ e si può portare dentro l'integrale interno, dove
$f_{X \mid Y}(x \mid y)\, f_Y(y) = f_{X,Y}(x, y)$:
\[ E[h(Y)] = \iint_{\R^2} g(x, y)\, f_{X,Y}(x, y) \dd x \dd y = E[g(X, Y)] . \]
\end{dimostrazione}

\begin{esempio}{La media in due tempi}
Con la densità $x + y$ sul quadrato: la media di $X$ sapendo $Y = y$ è
\[ E[X \mid Y = y] = \int_0^1 x\, \frac{x + y}{y + \frac12} \dd x = \frac{\frac13 + \frac{y}{2}}{y + \frac12} , \]
che vale $\frac23$ per $y = 0$ e $\frac59$ per $y = 1$: più $Y$ è grande, più la media di $X$ si abbassa. Mediando
rispetto a $Y$, il denominatore si semplifica con $f_Y(y) = y + \frac12$:
\[ E[X] = \int_0^1 E[X \mid Y = y]\, f_Y(y) \dd y = \int_0^1 \left( \tfrac13 + \tfrac{y}{2} \right) \dd y = \tfrac13 +
   \tfrac14 = \tfrac{7}{12} , \]
lo stesso valore che si ottiene direttamente: $\int_0^1 x \left( x + \frac12 \right) \dd x = \frac13 + \frac14 = \frac{7}{12}$.
\end{esempio}

Siano $(X, Y) \sim f_{X,Y}(x, y)$, con medie $(\mu_X, \mu_Y)$ e varianze $\sigma_X^2$, $\sigma_Y^2$. In analogia con il
caso discreto:
\begin{itemize}
  \item \textbf{covarianza}: $COV(X, Y) = E[(X - \mu_X)(Y - \mu_Y)] = E[XY] - \mu_X \mu_Y$;
  \item \textbf{coefficiente di correlazione}: $\rho_{XY} = \dfrac{COV(X, Y)}{\sigma_X\, \sigma_Y}$, con $\abs{\rho_{XY}} \leq 1$;
  \item \textbf{incorrelazione}: $COV(X, Y) = 0$;
  \item l'indipendenza statistica implica l'incorrelazione ma, in generale, non vale il viceversa.
\end{itemize}
Le dimostrazioni sono quelle del Capitolo~9, con gli integrali al posto delle somme.

\begin{esempio}{La covarianza della densità $x + y$}
Servono $E[XY]$ e le varianze. Con l'integrale doppio, a fette:
\[ E[XY] = \int_0^1 \int_0^1 x y (x + y) \dd y \dd x = \int_0^1 \int_0^1 \left( x^2 y + x y^2 \right) \dd y \dd x =
   \frac13 \cdot \frac12 + \frac12 \cdot \frac13 = \frac13 , \]
dove ogni termine è un prodotto di due integrali semplici. Quindi
$COV(X, Y) = \frac13 - \frac{7}{12} \cdot \frac{7}{12} = \frac{48 - 49}{144} = -\frac{1}{144}$. Con
$E[X^2] = \int_0^1 x^2 \left( x + \frac12 \right) \dd x = \frac14 + \frac16 = \frac{5}{12}$ si ha
$\sigma_X^2 = \frac{5}{12} - \frac{49}{144} = \frac{11}{144}$, e lo stesso per $Y$:
\[ \rho_{XY} = \frac{-1/144}{11/144} = -\frac{1}{11} \approx -0{,}09 . \]
Una correlazione negativa ma debole, coerente con le densità condizionate: quando $Y$ cresce, $X$ tende un poco a
calare.
\end{esempio}

% ============================================================
\section{La funzione di distribuzione congiunta}
% ============================================================

Come per una sola variabile, oltre alla densità si può usare una funzione che dà direttamente delle probabilità.

\dfn{CDF congiunta}{
  Data una coppia di variabili aleatorie $(X, Y)$, la \term{funzione di distribuzione di probabilità congiunta} è
  \[ F_{X,Y}(x, y) = P(X \leq x, Y \leq y) , \qquad (x, y) \in \R^2 . \]
  Descrive in modo completo il comportamento statistico congiunto delle due variabili.
}

$F_{X,Y}(x, y)$ è la probabilità che il punto $(X, Y)$ cada nel «quadrante» in basso a sinistra del punto $(x, y)$.

\begin{richiamo}{le derivate parziali}
Per una funzione di due variabili $F(x, y)$, la \term{derivata parziale} rispetto a $x$, $\frac{\partial F}{\partial x}$,
è la derivata fatta trattando $y$ come una costante (lo abbiamo già fatto nel Capitolo~7, derivando rispetto a $p$ con
$q$ fisso). Il simbolo $\partial$ (una «d» arrotondata) ricorda che le variabili sono più di una. La derivata
\term{mista} $\frac{\partial^2 F}{\partial x\, \partial y}$ si ottiene derivando prima rispetto a $y$ e poi il risultato
rispetto a $x$; per le funzioni regolari l'ordine non conta. Esempio: per $F = \frac{x^2 y + x y^2}{2}$,
$\frac{\partial F}{\partial y} = \frac{x^2 + 2xy}{2}$ e poi $\frac{\partial^2 F}{\partial x\, \partial y} = \frac{2x + 2y}{2}
= x + y$.
\end{richiamo}

Il legame con la densità è lo stesso del caso di una variabile, con una derivata per ogni variabile:
\[ F_{X,Y}(x, y) = \int_{-\infty}^{x} \int_{-\infty}^{y} f_{X,Y}(s, t) \dd t \dd s \qquad\Longleftrightarrow\qquad
   f_{X,Y}(x, y) = \frac{\partial^2}{\partial x\, \partial y} F_{X,Y}(x, y) . \]
Per esempio, per la densità $x + y$ sul quadrato, nei punti del quadrato
$F_{X,Y}(x, y) = \int_0^x \int_0^y (s + t) \dd t \dd s = \frac{x^2 y + x y^2}{2}$, e la sua derivata mista, calcolata nel
richiamo, ridà proprio $x + y$.

\teorema[proprietacdfc]{Proprietà della CDF congiunta}{
  \begin{enumerate}[label=\alph*)]
    \item $0 \leq F_{X,Y}(x, y) \leq 1$;
    \item fissato comunque $y_0$, $F_{X,Y}(x, y_0)$ è monotona non decrescente in $x$ (e continua da destra); fissato
          comunque $x_0$, $F_{X,Y}(x_0, y)$ è monotona non decrescente in $y$ (e continua da destra);
    \item $F_{X,Y}(-\infty, y) = F_{X,Y}(x, -\infty) = F_{X,Y}(-\infty, -\infty) = 0$ e
          $\lim_{x \to \infty, y \to \infty} F_{X,Y}(x, y) = 1$;
    \item le CDF marginali si ricavano dalla congiunta: $F_X(x) = F_{X,Y}(x, +\infty)$, $F_Y(y) = F_{X,Y}(+\infty, y)$;
    \item la probabilità dell'evento rettangolare $R = \{x_1 < X \leq x_2,\ y_1 < Y \leq y_2\}$ è
          \[ P(R) = F_{X,Y}(x_2, y_2) - F_{X,Y}(x_2, y_1) - F_{X,Y}(x_1, y_2) + F_{X,Y}(x_1, y_1) . \]
  \end{enumerate}
}

\begin{dimostrazione}
\passo{a), b), c)} $F_{X,Y}$ è una probabilità; allargando il quadrante (aumentando $x$ o $y$) l'evento si allarga e la
probabilità non diminuisce. Con $x = -\infty$ il quadrante è vuoto ($X \leq -\infty$ è impossibile); con entrambe le
coordinate a $+\infty$ è tutto il piano.
\passo{d)} $\{X \leq x, Y \leq +\infty\} = \{X \leq x\}$: la condizione su $Y$ è sempre vera.
\passo{e) Si toglie e si rimette} Il rettangolo si ottiene dal quadrante di $(x_2, y_2)$ togliendo il quadrante di
$(x_2, y_1)$ (la parte sotto) e quello di $(x_1, y_2)$ (la parte a sinistra). Così però il quadrante di $(x_1, y_1)$,
che sta in entrambi, è stato tolto due volte, e va rimesso una volta. È lo stesso ragionamento della probabilità
dell'unione (Capitolo~3).
\end{dimostrazione}

\begin{figure}[H]
\centering
\begin{tikzpicture}[x=1.15cm, y=1.15cm]
  \fill[mspiA!12] (-1.2,-1.2) rectangle (2.6,2.2);
  \fill[mspiB!25] (-1.2,-1.2) rectangle (2.6,0.6);
  \fill[mspiB!25] (-1.2,-1.2) rectangle (0.8,2.2);
  \fill[mspiB!50] (-1.2,-1.2) rectangle (0.8,0.6);
  \fill[evid] (0.8,0.6) rectangle (2.6,2.2);
  \draw[mspigray, dashed] (0.8,-1.3) -- (0.8,2.3); \draw[mspigray, dashed] (2.6,-1.3) -- (2.6,2.3);
  \draw[mspigray, dashed] (-1.3,0.6) -- (2.7,0.6); \draw[mspigray, dashed] (-1.3,2.2) -- (2.7,2.2);
  \node[font=\scriptsize, below] at (0.8,-1.3) {$x_1$}; \node[font=\scriptsize, below] at (2.6,-1.3) {$x_2$};
  \node[font=\scriptsize, left] at (-1.3,0.6) {$y_1$}; \node[font=\scriptsize, left] at (-1.3,2.2) {$y_2$};
  \node[font=\small] at (1.7,1.4) {$R$};
  \node[font=\scriptsize, align=center] at (-0.2,-0.3) {contato\\ 2 volte};
  \fill (0.8,0.6) circle (1.5pt); \fill (2.6,2.2) circle (1.5pt); \fill (0.8,2.2) circle (1.5pt); \fill (2.6,0.6) circle (1.5pt);
  \node[font=\small, anchor=west, text width=6.8cm] at (3.2,0.5)
    {Il quadrante di $(x_2, y_2)$ contiene tutto il disegno. Togliendo il quadrante di $(x_2, y_1)$ (la fascia in basso)
     e quello di $(x_1, y_2)$ (la fascia a sinistra) resta il rettangolo giallo $R$, ma l'angolo in basso a sinistra, il
     quadrante di $(x_1, y_1)$, è stato tolto due volte: lo si riaggiunge una volta.};
\end{tikzpicture}
\caption{La probabilità di un rettangolo dalla CDF congiunta.}
\end{figure}

\gbox{In sintesi}{mspiA!80!black}{
\begin{itemize}
  \item La \textbf{densità congiunta} $f_{X,Y}(x, y)$ è una probabilità per unità di area: $P((X, Y) \in C) = \iint_C
        f_{X,Y}$, il volume sotto la superficie. È $\geq 0$ e ha volume totale 1.
  \item L'\textbf{integrale doppio} si calcola a fette (Fubini): prima un integrale con l'altra variabile fissa, poi
        l'altro; per un prodotto $g(x)h(y)$ è il prodotto degli integrali.
  \item Marginali: $f_X = \int f_{X,Y} \dd y$. Indipendenza: $f_{X,Y} = f_X f_Y$ ($\iff F_{X,Y} = F_X F_Y$).
  \item Condizionata: $f_{X \mid Y}(x \mid y) = f_{X,Y}/f_Y$, ottenuta condizionando a $Y$ «vicino a $y$»; valgono
        probabilità totale, composta e Bayes per le densità.
  \item $E[g(X, Y)] = \iint g f_{X,Y}$, linearità, media condizionata $E[g] = E_Y[E_{X \mid Y}[g \mid Y]]$, covarianza e
        $\rho$ come nel discreto.
  \item \textbf{CDF congiunta} $F_{X,Y}(x, y) = P(X \leq x, Y \leq y)$, con $f_{X,Y} = \frac{\partial^2 F_{X,Y}}{\partial x
        \partial y}$; marginali $F_X(x) = F_{X,Y}(x, +\infty)$; un rettangolo si calcola con quattro valori della CDF.
\end{itemize}
}
